%%=============================================================================
%% Conclusie
%%=============================================================================

\chapter{Conclusie}%
\label{ch:conclusie}

Deze bachelorproef onderzocht in welke mate een Big Data-architectuur, gevoed door publiek beschikbare NBB-data, een meerwaarde biedt voor financiële risico-inschatting bij kredietverlening. Door de realisatie van een schaalbare Proof of Concept, waarbij een ETL-pijplijn data overbrengt naar Google BigQuery en verrijkt via Machine Learning-algoritmen, konden de limieten van traditionele, statische kredietanalyses overstegen worden.

% TODO: Vul deze sectie aan met de uiteindelijke resultaten van je PoC. Voorbeeld:
% Uit de resultaten van de implementatie blijkt dat de cloud-architectuur moeiteloos in staat is om complexe queries uit te voeren op miljoenen rijen aan Belgische bedrijfsdata. De exploratieve analyse toonde bovendien waardevolle sectorale trends aan die bij handmatige evaluaties onopgemerkt zouden blijven. Wat de predictieve analyses betreft, presteerden de getrainde ensemble-modellen naar verwachting beter dan de traditionele lineaire benaderingen...

De implementatie van dit platform demonstreert dat de financiële sector aanzienlijke voordelen kan halen uit het centraliseren en in-house bewerken van ruwe databronnen. In tegenstelling tot gesloten commerciële softwarepakketten, biedt een eigen datawarehouse de analist ongeziene flexibiliteit voor sectorbenchmarking en het vroegtijdig opsporen van financiële anomalieën. Deze proactieve benadering maakt het acceptatiebeleid rond zakelijke kredieten niet alleen sneller, maar bovenal objectiever.

Het onderzoek kent echter ook een aantal beperkingen. Ten eerste vereist het beheren van een dergelijk Big Data-platform aanzienlijke technische expertise op vlak van cloud-architectuur en data engineering, competenties die niet standaard aanwezig zijn bij klassieke risicoanalisten. Bovendien blijken historische jaarrekeningen niet altijd afdoende om plotse macro-economische schokken te voorspellen, gezien externe marktfactoren vaak niet vervat zitten in de NBB-statistieken.

Voor verder onderzoek is het ten zeerste aanbevolen om de opgezette dataset te verrijken met alternatieve of real-time databronnen, zoals transactiedata of nieuws-sentiment. Daarnaast biedt de integratie van geautomatiseerde waarschuwingssystemen (alerting pipelines) rechtstreeks in de interface van de analist een interessant traject voor toekomstige ontwikkeling.